\section{Diseño del controlador y evaluación MIL}

\pendiente{Toda esta sección usa el modelo del diseño ($c = 5/7$, retardo de
\SI{75}{\milli\second}), porque el diseño se hizo antes de identificar la planta. Cuando se
cierre el rediseño con el modelo nominal de la Sección~II-D, decidir si esta sección presenta
los dos diseños (original y rediseño) bajo las mismas condiciones o solo el definitivo.}

\subsection{Estructura y sintonización}

Cada eje es un doble integrador con retardo, así que la estructura se eligió a partir de esa
dinámica y no del Proyecto~1. La acción derivativa es obligatoria para tener margen de fase. En
vez de derivar la posición medida, el controlador usa la velocidad $\vh$ que ya estima el filtro
de Kalman de la PC, con lo que no se agrega otro filtro ni más atraso. La acción integral hace
falta para corregir sesgos como el desnivel del plato o el neutro de los servos: sin ella, cada
grado de desnivel deja un error de unos \SI{29}{\milli\metre}. Pero un PID clásico no sirve en
esta planta, porque al arrancar con la bola lejos de la referencia el integrador se carga durante
todo el recorrido. Por eso se usó \emph{integración condicional}: el integrador solo acumula
después de cuatro cuadros seguidos con la bola casi quieta, $|\vh| < \SI{25}{\milli\metre\per\second}$,
que es cuando el error que queda es un sesgo y no un transitorio. La Tabla~\ref{tab:estructuras}
compara las tres estructuras con las mismas $\Kp$ y $\Kd$ en el simulador no lineal.

\begin{table}[!t]
\caption{Comparación de estructuras en el simulador no lineal, con las mismas $\Kp$ y $\Kd$.
Regulación desde \SI{100}{\milli\metre} y desnivel de \SI{1}{\degree}. Fuente:
\texttt{etapaC/resultados/C1\_consola.txt}.}
\label{tab:estructuras}
\centering
\footnotesize
\setlength{\tabcolsep}{3pt}
\begin{tabular}{@{}lS[table-format=1.2]S[table-format=2.1]S[table-format=1.2]S[table-format=2.2]@{}}
\toprule
 & \multicolumn{3}{c}{\textbf{Regulación \SI{100}{\milli\metre}}} & {\textbf{Desnivel}} \\
\cmidrule(lr){2-4}\cmidrule(l){5-5}
\textbf{Estructura} & {$t_s$ (\si{\second})} & {$M_p$ (\si{\percent})} & {$e_{ss}$ (\si{\milli\metre})} & {$e_{ss}$ (\si{\milli\metre})} \\
\midrule
PD                       & 1.40 & 2.3  & 0.28 & 30.52 \\
PID clásico              & 3.93 & 38.8 & 2.07 & 0.14 \\
PID integr. condicional  & 1.30 & 4.4  & 0.18 & 0.27 \\
\bottomrule
\end{tabular}
\end{table}

$\Kp$ y $\Kd$ se calcularon por asignación de polos. El retardo del lazo se representó como un
atraso equivalente $t_e = t_d + \Ts/2 + 2/\omega_{n,\mathrm{srv}}$, con $t_d = \SI{75}{\milli\second}$
y $\omega_{n,\mathrm{srv}} = \SI{50}{\radian\per\second}$ supuesto para el servo, lo que da
$t_e = \SI{0.132}{\second}$. Con el modelo de diseño $-c\,g/(s^{2}(t_e s + 1))$, el polinomio
característico $t_e s^{3} + s^{2} + c\,g\,\Kd\,s + c\,g\,\Kp$ se igualó a
$t_e(s + a\omega)(s^{2} + 2\zeta\omega s + \omega^{2})$ con $\zeta = 0{,}7$ y $a = 1$. El
coeficiente de $s^{2}$ fija $\omega = 1/(t_e(a + 2\zeta))$: es el retardo el que limita la
velocidad del lazo. La integral se ajustó con $\Ti = \SI{1.5}{\second}$, con el error que entra
al integrador recortado a $\pm\SI{30}{\milli\metre}$. Los parámetros resultantes y los márgenes
del lazo discreto, con la planta ZOH, el retardo, el Kalman y el controlador lineal, se
resumen en la Tabla~\ref{tab:controlador}. En la parte lineal, con $C(s)$ en \si{\radian\per\milli\metre}, el controlador es

\begin{equation}
C(s) = \frac{4{,}52\times10^{-4}\,s^{2} + 5{,}95\times10^{-4}\,s + 3{,}97\times10^{-4}}{s}
.
\label{eq:cs}
\end{equation}

\begin{table}[!t]
\caption{Parámetros del controlador y márgenes del lazo discreto nominal. Fuente:
\texttt{etapaC/resultados/C1\_consola.txt} (MATLAB R2026a).}
\label{tab:controlador}
\centering
\footnotesize
\begin{tabular}{@{}lr@{}}
\toprule
\textbf{Magnitud} & \textbf{Valor} \\
\midrule
$\Kp$ & \SI{0.03412}{\degree\per\milli\metre} \\
$\Kd$ & \SI{0.02588}{\degree\per(\milli\metre\per\second)} \\
$\Ki$ (continua) & \SI{0.02275}{\degree\per(\milli\metre\second)} \\
$\Ki\Ts$ (incremento por muestra) & \SI{7.58e-4}{\degree\per\milli\metre} \\
Saturación $\thlim$ & \SI{10}{\degree} \\
Kalman ($\sigma_a = \SI{500}{\milli\metre\per\second\squared}$, $\sigma_m = \SI{1}{\milli\metre}$) & $K = [\num{0.647};\ \SI{9.90}{\per\second}]$ \\
\midrule
Frecuencia de cruce $\omega_c$ & \SI{3.30}{\radian\per\second} \\
Margen de fase & \SI{40.9}{\degree} \\
Margen de ganancia superior / inferior & \SI{+11.8}{} / \SI{-12.6}{\deci\bel} \\
Margen de retardo & \SI{216}{\milli\second} \\
PM con retardo de 90 / \SI{120}{\milli\second} & \SI{38.0}{} / \SI{32.3}{\degree} \\
\bottomrule
\end{tabular}
\end{table}

El margen de ganancia inferior aparece porque el lazo es condicionalmente estable: si la ganancia
de la planta baja más de \SI{12.6}{\deci\bel}, el doble integrador se desestabiliza. Este dato
importa en la Sección~VII, porque la identificación dio una ganancia efectiva bastante menor que
la nominal.

La ley discreta que se ejecuta en cada cuadro, por eje, es

\begin{equation}
\begin{aligned}
\varepsilon_k &= \xh_k - r_k, \qquad u_k = \Kp\,\varepsilon_k + I_k + \Kd\,\vh_k,\\
\theta_{\mathrm{cmd},k} &= \operatorname{sat}(u_k, \pm\thlim),\\
n_k &= \begin{cases} n_{k-1} + 1 & |\vh_k| < \SI{25}{\milli\metre\per\second}\\ 0 & \text{en otro caso}\end{cases}\\
I_{k+1} &= I_k + \Ki\Ts\,\operatorname{sat}(\varepsilon_k, \pm\SI{30}{\milli\metre}),\\
&\quad\text{solo si } n_k \ge 4 \text{ y no hay \emph{clamping}},
\end{aligned}
\label{eq:ley}
\end{equation}

donde el \emph{clamping} congela el integrador cuando la salida está saturada y el error tiene el
mismo signo que $u_k$ \cite{mw_antiwindup}.

\pendiente{Si el enunciado pide un diseño alternativo por otro método, presentarlo aquí. El
rediseño de la Sección~VII es la misma estructura con otra sintonía, no una estrategia distinta.
Candidatos reales: realimentación de estados con observador (LQR/LQG) o un compensador de
adelanto--atraso por respuesta en frecuencia.}

\subsection{Evaluación Model-in-the-Loop}

El MIL simula la cadena completa con el modelo no lineal de dos ejes: la cámara toma muestras
cada $\Ts$ con ruido de \SI{1}{\milli\metre}, el dato llega a la ESP32 con el retardo, pasa por el
Kalman y el controlador, la inclinación se convierte en comando con la tabla servo--plato medida,
y el servo, de segundo orden con $\omega_n = \SI{50}{\radian\per\second}$, velocidad de
\SI{600}{\degree\per\second} y carrera de $\pm\SI{45}{\degree}$, mueve el plato. Todos los casos
comparten planta, $\Ts$, retardo, ruido (con la misma semilla), integrador numérico y duración,
salvo lo que cada caso cambia a propósito. La Fig.~\ref{fig:mil} muestra el ensayo de requisitos
y el de saturación, y la Tabla~\ref{tab:mil} las métricas.

\begin{figure*}[!t]
\centering
\begin{tikzpicture}
\begin{groupplot}[group style={group size=2 by 2, horizontal sep=1.3cm, vertical sep=0.9cm},
                  estilo, ancho, xmin=0, xmax=6]
\nextgroupplot[ylabel={Posición (\si{\milli\metre})}, xticklabels={},
               title={\scriptsize D1: regulación desde \SI{100}{\milli\metre}}, title style={yshift=-2pt},
               legend pos=north east]
  \addplot[cAzul] table[col sep=comma, x=t_cap_s, y=xh_x_mm] {\G/D_caso_D1_requisitos.csv};
  \addplot[cVerde, dashed] table[col sep=comma, x=t_cap_s, y=x_plato_mm] {\G/D_caso_D1_requisitos.csv};
  \addplot[cGris, densely dotted] table[col sep=comma, x=t_cap_s, y=r_x_mm] {\G/D_caso_D1_requisitos.csv};
  \legend{$\xh$, $x$ (plato), $r$}
\nextgroupplot[ylabel={Posición (\si{\milli\metre})}, xticklabels={},
               title={\scriptsize D2: empujón con saturación y anti-windup}, title style={yshift=-2pt},
               legend pos=north east]
  \addplot[cAzul] table[col sep=comma, x=t_cap_s, y=xh_x_mm] {\G/D_caso_D2_sat_con_AW.csv};
  \addplot[cVerde, dashed] table[col sep=comma, x=t_cap_s, y=x_plato_mm] {\G/D_caso_D2_sat_con_AW.csv};
  \addplot[cRojo, densely dotted, domain=0:6, samples=2] {127};
  \legend{$\xh$, $x$ (plato), Borde}
\nextgroupplot[ylabel={$\thc$ (\si{\degree})}, xlabel={Tiempo (\si{\second})}]
  \addplot[cNaranja] table[col sep=comma, x=t_cap_s, y=thcmd_x_deg] {\G/D_caso_D1_requisitos.csv};
\nextgroupplot[ylabel={$\thc$ (\si{\degree})}, xlabel={Tiempo (\si{\second})}]
  \addplot[cNaranja] table[col sep=comma, x=t_cap_s, y=thcmd_x_deg] {\G/D_caso_D2_sat_con_AW.csv};
  \addplot[cRojo, densely dotted, domain=0:6, samples=2] {10};
\end{groupplot}
\end{tikzpicture}
\caption{Evaluación MIL con el actuador medido. Izquierda: ensayo de requisitos. Derecha:
empujón de \SI{420}{\milli\metre\per\second} con \SI{1}{\degree} de desnivel, con saturación a
$\thlim = \SI{10}{\degree}$ y anti-windup.}
\label{fig:mil}
\end{figure*}

\begin{table}[!t]
\caption{Métricas MIL (eje X). Fuente: \texttt{etapaD/resultados/D1\_metricas.csv}
(MATLAB R2026a).}
\label{tab:mil}
\centering
\footnotesize
\setlength{\tabcolsep}{2.5pt}
\begin{tabular}{@{}lccccc@{}}
\toprule
\textbf{Caso} & $t_s$ (\si{\second}) & $M_p$ (\si{\percent}) & $e_{ss}$ (\si{\milli\metre}) & $\max|x|$ (\si{\milli\metre}) & $t_{\mathrm{sat}}$ (\si{\second}) \\
\midrule
D1 requisitos                 & 1,30 & 4,4  & 0,18 & 102,7 & 0 \\
D2 sin saturación             & ---  & ---  & 1,48 & 127,1 & 0 \\
D2 con sat., sin AW           & ---  & ---  & 4,49 & 133,0 & 0,30 \\
D2 con sat., con AW           & ---  & ---  & 4,49 & 133,0 & 0,30 \\
D2b PID clásico, sin AW       & ---  & ---  & 2,97 & 133,0 & 0,33 \\
D2b PID clásico, con AW       & ---  & ---  & 2,40 & 133,0 & 0,30 \\
D3 desnivel de \SI{1}{\degree} & ---  & ---  & 0,27 & 28,8  & 0 \\
\midrule
D5 retardo de \SI{120}{\milli\second} & 3,03 & 5,0  & 0,14 & 102,6 & 0 \\
D5 ruido de \SI{2}{\milli\metre}      & 4,90 & 6,5  & 0,96 & 102,9 & 0 \\
D5 firmware sin tabla                 & 4,73 & 23,8 & 0,78 & 80,2  & 0 \\
\bottomrule
\end{tabular}
\end{table}

En el ensayo de requisitos el controlador cumple R1 a R4: $t_s = \SI{1.30}{\second}$,
$M_p = \SI{4.4}{\percent}$ y $e_{ss} = \SI{0.18}{\milli\metre}$, con un pico de
\SI{3.8}{\degree} de inclinación, lejos de la saturación. La saturación solo aparece con
perturbaciones que llevan la bola casi al borde. Con el empujón de D2, el plato satura
\SI{0.30}{\second} y la bola llega a \SI{133}{\milli\metre}, más allá del borde: con
$\thlim = \SI{10}{\degree}$, el actuador real no alcanza a frenar esa perturbación, y eso es una
limitación del mecanismo y no de la sintonía. Con integración condicional, las versiones con y sin
anti-windup dan exactamente lo mismo, porque durante la saturación la bola se mueve rápido y el
integrador ya está detenido por la condición de velocidad. Con un PID clásico, en cambio, el
anti-windup sí actúa y baja el error final de \SI{2.97}{} a \SI{2.40}{\milli\metre}.

Los casos de robustez muestran dónde está el riesgo: con \SI{120}{\milli\second} de retardo o con
\SI{2}{\milli\metre} de ruido, $t_s$ pasa de \SI{3}{\second}, y si el firmware convirtiera la
inclinación en comando con la pendiente del diseño en vez de la tabla medida, el sobreimpulso
subiría a \SI{23.8}{\percent}. Por eso el firmware usa la tabla.

\pendiente{Verificar R4 en D2 de forma explícita en el texto final (no se cumple con ese
empujón). Agregar el caso D4 (dos ejes) si cabe en la extensión.}
